\subsection{Asymptotic accuracy of TDS -- additional details and full theorem statement}\label{sec:theorem_statement_and_proof}
In this section we (i) characterize sufficient conditions on the model and twisting functions under which TDS provides arbitrarily accurate estimates as the number of particles is increased and 
(ii) discuss when these conditions will hold in practice for the twisting function $\pTwist(y\mid \xt)$ introduced in \Cref{sec:method}.

We begin with a theorem providing sufficient conditions for asymptotic accuracy of TDS.
\begin{theorem}\label{thm:twisted_SMC_complete}
Let $\pmodel(\dt{x}{0:T})$ be a diffusion generative model
(defined by \cref{eqn:diffusion_marginal,eq:unconditional-diffusion-transition})
with 
$$
\pmodel(\xt \mid \xtpo) = \mathcal{N}\left(\xt \mid \xtpo + \sigma_{t+1}^2 \scoremodel(\xtpo), \sigma_{t+1}^2 I
\right),$$
with variances $\sigma_1^2, \dots, \sigma_T^2.$ 
Let $\pTwist(y \mid \xt)$ be twisting functions,
and 
$$
\proposal_t(\xt \mid \xtpo) = 
\mathcal{N}\left(
\xt \mid \xtpo + \sigma_{t+1}^2 [\scoremodel(\xtpo) + \nabla_{\xtpo}\log \pTwist(y\mid \xtpo)], 
\tilde \sigma_{t+1}^2 \right)
$$
be proposals distributions for $t=0,\dots, T-1,$
and let $\mathbb{P}_{K}=\sum_{k=1}^K \dt{w_k}{0}\delta_{\dt{x_k}{0}}$ for weighted particles $\{(\dt{x_k}{0}, \dt{w_k}{0})\}_{k=1}^K$ returned by \Cref{alg:TDS} with $K$ particles.
Assume
\begin{enumerate}[label=(\alph*)]
    \item{the final twisting function is the likelihood, $\pTwist(y\mid \xzero)=\plikelihood{y}{\xzero},$} \label{item:final_twist}
    \item{the first twisting function $\pTwist(y \mid \xT),$ and the ratios of subsequent twisting functions $\pTwist(y \mid \xt) / \pTwist(y \mid \xtpo) $ are positive and bounded,}\label{item:twist_assumption}
    %\item{each $\nabla_{\xt} \log \pTwist(y\mid \xt)$ with $t>0$ is continuous and has bounded gradients, and}\label{item:twist_gradient}
    \item{each $\log \pTwist(y\mid \xt)$ with $t>0$ is continuous and has bounded gradients in $\xt$, and}\label{item:twist_gradient}
    \item{the proposal variances are larger than the model variances, i.e.\ for each $t,\ \tilde \sigma_t^2 >\sigma_t^2.$} \label{item:proposal_var}
\end{enumerate}
Then $\mathbb{P}_{K}$ converges setwise to $\pmodel(\xzero\mid y)$ with probability one,
that is 
for every set $A,$ 
$\lim_{K\rightarrow\infty} \mathbb{P}_{K}(A) = \int_A \pmodel(\xzero\mid y) d\xzero.$
\end{theorem}

The assumptions of \Cref{thm:twisted_SMC_complete} are readily satisfied in common applications.
\begin{itemize}
\item {Assumption \ref{item:final_twist} may be satisfied by construction by choosing $\pTwist(y\mid \xzero)=\plikelihood{y}{\xzero}$
by defining $\denoiser(\xzero, t=0)=\xzero.$
}
\item{Assumption \ref{item:twist_assumption} is satisfied if (i) $\plikelihood{y}{x}$ is smooth in $x$ and everywhere positive and (ii) $\pTwist(y\mid \xt) = \plikelihood{y}{\denoiser(\xt)}$ where $\denoiser(\xt)$ takes values in some compact domain.
An alternative sufficient condition for Assumption \ref{item:twist_assumption} is for $\plikelihood{y}{x}$ to be positive and bounded away from zero; 
this latter condition will hold when, for example, $\plikelihood{y}{x}$ is a classifier fit with regularization.}
\item{Assumption \ref{item:twist_gradient} is the strongest assumption.
It will be satisfied, for example, if (i) $\pTwist(y\mid \xt) = \plikelihood{y}{\denoiser(\xt)},$ and (ii) $\plikelihood{y}{x}$ and $\denoiser(\xt)$ are smooth in $x$ and $\xt,$ with uniformly bounded gradients.
While smoothness of $\denoiser(\cdot, t)$ can be encouraged by the use of skip-connections and regularization,
$\denoiser(\cdot, t)$ may present sharp transitions,
particularly for $t$ close to zero.}
\item{Assumption \ref{item:proposal_var}, that the proposal variances satisfy $\tilde \sigma_{t}^2 > \sigma_t^2$ is likely not needed for the result to hold.
However, this assumption permits usage of existing SMC theoretical results in the proof;
in practice, our experiments use $\tilde \sigma_t^2 = \sigma_t^2,$
but alternatively the assumption could be met by inflating each $\tilde \sigma_t^2$ by some arbitrarily small $\delta$ without markedly impacting the behavior of the sampler.}
\end{itemize}

%\paragraph{\BLT{Cut this paragraph for NeurIPS appendix submission} Inpainting and inpaiting with degrees of freedom.}

%\begin{cor}
%The inpainting and DOF cases are consistent for their conditionals.
%\end{cor}
%\begin{proof}
%We can see that early truncation at step $t=1$ forms a special case.
%By reindexing $t$ we find we are consistent for $\pmodel(\dt{x}{1} \mid y).$
%Then those final samplers sample $\xzero \sim \pmodel(\xzero \mid \dt{x}{1}, y)$ in the last step.
%So by the law of total probability and chain rule of probability they give exact samples in the limit as desired.
%\end{proof}
%
%Why is bounded assumption reasonable in the inpainting case?
%It's because each $\sigma_t^2>0$ and so with 
%$\pTwist(y\mid \xt; \mask) = \mathcal{N}(y; \denoiser(\xt)_\mask, \bar \sigma_t^2)=p(y\mid \dt{x}{1})$
%with $\bar \sigma_t^2 = \sum_{t^\prime=1}^t\sigma_{t^\prime}^2,$
%the twisting variance are a decreasing sequence.
%And since the denoiser is also bounded and each $\bar \sigma_t^2>0$, there is some $0 <c <C< \infty$ such that 
%$c<\pTwist(y\mid \xt) <C$ and so the ratio is bounded.

\paragraph{Proof of \Cref{thm:twisted_SMC_complete}:}
\Cref{thm:twisted_SMC_complete} characterizes a set of conditions under which SMC algorithms converge.
We restate this result below in our own notation.
\begin{theorem}[\citet{chopin2020introduction} -- Proposition 11.4]\label{thm:chopin_theorem}
Let $\{(\dt{x_k}{0}, \dt{w_k}{0})\}_{k=1}^K$ be the particles and weights returned at the last iteration of a sequential Monte Carlo algorithm with $K$ particles using multinomial resampling.
If each weighting function $\weight_t(\xt, \xtpo)$ is positive and bounded, 
then for every bounded, $\target_0$-measurable function $\phi$ of $\xt$
$$
\lim_{K\rightarrow \infty} \sum_{k=1}^K \dt{w_k}{0} \phi(\dt{x_k}{0})
=
\int \phi(\xzero)\target_0(\xzero)d\xzero.
$$
with probability one.
\end{theorem}


An immediate consequence of \Cref{thm:chopin_theorem} is the setwise convergence of the discrete measures, $\hat P_K = \sum_{k=1}^K \dt{w_k}{0} \delta_{\dt{x_k}{0}}.$
This can be seen by taking for each $\phi(x) = \I[x \in A]$ for any $\target_0$-measurable set $A$.
The theorem applies both in the Euclidean setting, where each $\dt{x_k}{t}\in \R^D,$ as well as the Riemannian setting.

We now proceed to prove \Cref{thm:twisted_SMC_complete}.
\begin{proof}
To prove the theorem we show
(i) the $\xzero$ marginal final target $\target_0$ is $\pmodel(\xzero \mid y)$ and then
(ii) $\mathbb{P}_K$ converges setwise to $\target_0.$

We first show (i) by manipulating $\target_0$ in \Cref{eq:final-target} to obtain $\pmodel(\dt{x}{0:T} \mid y).$
From \Cref{eq:final-target} we first have
\begin{align}
\begin{split}
   \target_0(\dt{x}{0:T}) &= \frac{1}{\mathcal{L}_0} 
        \left[ \proposal(\dt{x}{T})  \prod_{t=0}^{T-1} \proposal_{t}(\dt{x}{t}\mid  \dt{x}{t+1})  \right] 
        \left[ \weight_T(\dt{x}{T})
        \prod_{t=0}^{T-1} \weight_{t^\prime}(\dt{x}{t}, \dt{x}{t+1}) \right]   \\
    &\mathrm{Substitute\ in\ weights\ from\ \cref{eq:diffusion-twisted-weigthing-functions}.} \\
    &= \frac{1}{\mathcal{L}_0} 
        \left[ p(\dt{x}{T}) \prod_{t=0}^{T-1} r_t (\xt \mid \xtpo)
        \right] 
        \left[ \pTwist(y\mid \dt{x}{T})
        \prod_{t=0}^{T-1} 
        \frac{\pmodel(\xt \mid \xtpo)\pTwist(y\mid \xt)}{\pTwist(y \mid \xtpo)r_t(\xt \mid \xtpo)}
        \right] \\
    &\mathrm{Rearrange\ }\pmodel\mathrm{\ and \ } r_t\ \mathrm{\ terms, \ and\ cancel\ out\ }r_t \mathrm{\ terms.} \\
    &= \frac{1}{\mathcal{L}_0} 
        \left[ p(\dt{x}{T}) \prod_{t=0}^{T-1} \pmodel(\xt \mid \xtpo)
        \right] 
        \left[ \pTwist(y\mid \dt{x}{T})
        \prod_{t =0}^{T-1} 
        \frac{\cancel{r_t(\xt \mid \xtpo)}\pTwist(y\mid \xt)}{\pTwist(y \mid \xtpo)\cancel{r_t(\xt \mid \xtpo)}}
        \right] \\
    &\mathrm{Collapse\ }\pmodel\mathrm{\ terms.}\\ %,\ rearrange\ and\ cancel\ }\pTwist \mathrm{\ terms.} \\
    &= \frac{1}{\mathcal{L}_0} 
        \pmodel(\dt{x}{0:T})
        \left[ 
        \prod_{t=0}^{T-1} 
        \frac{\pTwist(y\mid \xt)}{\pTwist(y \mid \xtpo)}\right]\pTwist(y \mid \dt{x}{T})\\
    &\mathrm{Cancel \ out \ } \pTwist \mathrm{\ terms.}\\
    &= \frac{1}{\mathcal{L}_0} 
        \pmodel(\dt{x}{0:T}) \pTwist (y \mid \xzero) \\ 
    &\mathrm{Recognize\ }\pTwist(y\mid \dt{x}{0})=\plikelihood{y}{\xzero}\mathrm{\ by \ Assm. \ref{item:final_twist}\ \ and\  apply\  Bayes'\  rule\  with\  }\mathcal{L}_0=\pmodel(y). \\
    &= \pmodel(\dt{x}{0:T}\mid y).
\end{split}
\end{align}
The final line reveals that once we marginalize out $\dt{x}{1:T}$
we obtain $\target_0(\xzero)=p(\xzero \mid y)$ as desired.

We next show that $\mathbb{P}_K$ converges to $\target_0$ with probability one by applying \Cref{thm:chopin_theorem}.
To apply \Cref{thm:chopin_theorem} it is sufficient to show that the weights at each step are upper bounded, as they are defined through (ratios of) probabilities and hence are positive. 
Since there are a finite number of steps $T,$
it is enough to show that each $w_t$ is bounded.
The inital weight is the initial twisting function, $w_T(\xT)=\pTwist(y\mid \xT),$ which is bounded by Assumption \ref{item:twist_assumption}.
So we proceed to intermediate weights.

To show that the weighting functions at subsequent steps are bounded,
we decompose the log-weighting functions as 
$$
\log \weight_t(\xt, \xtpo) = \log \frac{\pTwist(y\mid \xt)}{\pTwist(y\mid \xtpo)} + \log \frac{\pmodel(\xt\mid \xtpo)}{r_t(\xt \mid \xtpo)} ,
$$
and show independently that $\log \pTwist(y\mid \xt)/\pTwist(y\mid \xtpo)$ and
$\log \pmodel(\xt\mid \xtpo)/r_t(\xt \mid \xtpo)$ are bounded.
The first term $\log \pTwist(y\mid \xt)/\pTwist(y\mid \xtpo)$ is again bounded by Assumption \ref{item:twist_assumption},
and we proceed to the second.

That $\log \pmodel(\xt\mid \xtpo)/r_t(\xt \mid \xtpo)$ is bounded follows from Assumptions \ref{item:twist_gradient} and \ref{item:proposal_var}.
First write
$$
%\pmodel(\xt \mid \xtpo) = \mathcal{N}\left(\xt \mid \xtpo + \sigma_{t+1}^2 \scoremodel(\xtpo), \sigma_{t+1}^2 I
\pmodel(\xt \mid \xtpo) = \mathcal{N}\left(\xt \mid \hat \mu, \sigma_{t+1}^2 I
\right)$$
with $\hat \mu=\xtpo + \sigma_{t+1}^2 \scoremodel(\xtpo),$
and 
$$
\proposal_t(\xt \mid \xtpo) = 
\mathcal{N}\left(
\xt \mid  \hat \mu_\psi, \tilde \sigma_{t+1}^2 I \right),
$$
for $\hat \mu_\psi = \hat \mu + \sigma_{t+1}^2 \nabla_{\xtpo}\log \pTwist(y\mid \xtpo)$.
The log-ratio then simplifies as 
\begin{align}
\log \frac{\pmodel(\xt\mid \xtpo)}{r_t(\xt \mid \xtpo)} &= 
\log \frac{| 2\pi \sigma_{t+1}^2 I|^{-1/2} \exp\{-(2\sigma_{t+1}^2)^{-1}\|\hat \mu - \xt \|^2 \}}
{ | 2\pi \tilde \sigma_{t+1}^2 I|^{-1/2} \exp\{-(2\tilde \sigma_{t+1}^2)^{-1}\|\hat \mu_\psi- \xt \|^2 \}} \\
&\textrm{Rearrange \ and\  let  \ } C=\log | 2\pi \sigma_{t+1}^2 I|^{-1/2} /  | 2\pi \tilde \sigma_{t+1}^2 I|^{-1/2}\\
&= \frac{-1}{2}\left[ 
\sigma_{t+1}^{-2}\|\hat \mu - \xt \|^2    -
\tilde \sigma_{t+1}^{-2}\|\hat \mu_\psi - \xt \|^2 \right]  +C \\
&\textrm{Expand and rearrange \ } \|\hat \mu_\psi - \xt \|^2 = \| \hat \mu  - \xt \|^2 + 2 \langle \hat\mu_\psi - \hat \mu, \hat \mu - \xt  \rangle + \|\hat \mu_\psi - \hat \mu\|^2 \\
&= \frac{-1}{2}\left[ ( \sigma_{t+1}^{-2} - \tilde \sigma_{t+1}^{-2})\|\hat \mu - \xt \|^2  
- 2 \tilde \sigma_{t+1}^{-2}\langle \hat\mu_\psi - \hat \mu, \hat \mu - \xt  \rangle - \tilde \sigma_{t+1}^2 \|\hat \mu_\psi - \hat \mu\|^2 
\right]+ C\\
&\textrm{Let \ } C^\prime {=} C - \frac{1}{2}\tilde \sigma_{t+1}^2 \|\hat \mu_\psi - \hat \mu\|^2 \textrm{\ and rearrange.  Note that $\|\hat \mu_\psi {-} \hat \mu\|^2{<}\infty$ by Assm. \ref{item:twist_gradient}.}\\
&= \frac{-1}{2}(\sigma_{t+1}^{-2} - \tilde \sigma_{t+1}^{-2})\|\hat \mu - \xt \|^2 +
\tilde \sigma_{t+1}^{-2}\langle \hat\mu_\psi - \hat \mu, \hat \mu - \xt  \rangle + C^\prime\\
&\textrm{Apply Cauchy-Schwarz}\\
&\le  \frac{-1}{2} (\sigma_{t+1}^{-2} - \tilde \sigma_{t+1}^{-2})\|\hat \mu - \xt \|^2
+ \tilde \sigma_{t+1}^{-2}\| \hat\mu_\psi - \hat \mu\|\cdot \| \hat \mu - \xt  \| + C^\prime\\
&\textrm{Upper-bounding using that $\mathrm{max}_{x}\frac{-a}{2}x^2 + bx =\frac{b^2}{2a}$ for $a = \sigma_{t+1}^{-2} - \tilde \sigma_{t+1}^{-2}>0$ by Assm. (d).}\\
&\le  \frac{1}{2} \frac{(\tilde \sigma_{t+1}^{-2}\| \hat\mu_\psi - \hat \mu\|)^2}{\sigma_{t+1}^{-2} - \tilde \sigma_{t+1}^{-2}} + C^\prime \\
&= \frac{\tilde \sigma_{t+1}^{-4}}{2( \sigma_{t+1}^{-2} - \tilde \sigma_{t+1}^{-2})}\| \hat\mu_\psi - \hat \mu\|^2  + C^\prime \\
& \textrm{Note that } \hat \mu_\psi = \hat \mu + \sigma_{t+1}^2 \nabla_{\xtpo}\log \pTwist(y\mid \xtpo). \\ 
&= \frac{\tilde \sigma_{t+1}^{-4}}{2( \sigma_{t+1}^{-2} - \tilde \sigma_{t+1}^{-2})}\sigma_{t+1}^4 \| \nabla_{\xtpo} \log \pTwist(y \mid \xtpo) \|^2  + C^\prime \\
&\le C^{\prime\prime}.
\end{align}
The final line follows from Assumption \ref{item:twist_gradient}, that the gradients of the twisting functions are bounded.
The above derivation therefore provides that each $\weight_t$ is bounded, concluding the proof.
\end{proof}